\documentclass[11pt,a4paper]{article}

\usepackage[utf8]{inputenc}
\usepackage[T1]{fontenc}
\usepackage[spanish]{babel}
\usepackage[margin=2.5cm]{geometry}
\usepackage{amsmath,amssymb}
\usepackage{booktabs,tabularx,array}
\usepackage{enumitem}
\usepackage{xcolor}
\usepackage{listings}
\usepackage[hidelinks]{hyperref}

\lstset{
  basicstyle=\ttfamily\small,
  frame=single,
  breaklines=true,
  columns=fullflexible,
  keepspaces=true,
  backgroundcolor=\color{black!4},
  rulecolor=\color{black!25}
}
\newcolumntype{Y}{>{\raggedright\arraybackslash}X}
\newcommand{\code}[1]{\texttt{#1}}

\title{Simple Oficios: servicio \code{mapas}\\[2pt]
\large Directorio de talleres, ranking de mecánicos, agentes y decisiones de diseño}
\author{Equipo Simple Oficios}
\date{29 de octubre de 2026}

\begin{document}
\maketitle

\begin{abstract}
Este documento describe la parte del proyecto que vive en la carpeta \code{mapas/}: un servicio
de API que ofrece un directorio de talleres mecánicos filtrable por estado, un ranking de
mecánicos basado en un orden lexicográfico, los agentes de inteligencia artificial que conducen
la conversación con el cliente, y el modelo de decisión con el que se eligió cómo construirlo.
Se indican también los límites de lo ya verificado y lo que sigue pendiente.
\end{abstract}

\tableofcontents

%------------------------------------------------------------------------------
\section{Alcance y premisas}

Simple Oficios es un marketplace que conecta clientes con mecánicos confiables. Esta parte del
proyecto cubre únicamente la búsqueda y recomendación de talleres y los agentes. Las premisas
vigentes son:

\begin{itemize}[nosep]
  \item Piloto en la Ciudad de México (CDMX); el filtro por estado aplica a talleres y mecánicos.
  \item Los \emph{leads} para los mecánicos son gratuitos: no hay pasarela de pagos.
  \item La base de datos (Supabase) y sus pruebas las lleva otra persona del equipo; aquí solo
        existe una propuesta de esquema, no aplicada.
  \item El despliegue es en Vercel, como un servicio dentro de un proyecto con varios servicios.
  \item Los datos geográficos son abiertos: el directorio DENUE y la Red Nacional de Caminos, ambos de INEGI.
\end{itemize}

%------------------------------------------------------------------------------
\section{Arquitectura}

El servicio es una aplicación FastAPI. Vercel entrega a cada servicio la ruta original de la
petición, por lo que todas las rutas llevan el prefijo \code{/api/mapas}.

\begin{lstlisting}
simple_oficios/
  vercel.json              servicio "mapas" y su rewrite
  mapas/
    main.py                API FastAPI bajo /api/mapas
    config.py              modelo de lenguaje (cualquier API compatible con OpenAI)
    agents/                diagnostico, matching, agenda y orquestador (Strands)
    tools/                 estados, SCIAN, DENUE, ranking, cargadores de datos
    decisiones/            modelo lexicografico y DECISIONES.md
    tests/                 26 pruebas sin red ni claves
    data/prod/             base minima de talleres de CDMX (sin telefonos)
    docs/                  propuesta de esquema para Supabase (no aplicada)
\end{lstlisting}

\begin{table}[h]
\centering
\begin{tabularx}{\textwidth}{@{}l l Y@{}}
\toprule
Método & Ruta & Descripción \\
\midrule
GET  & \code{/api/mapas/health}    & Estado del servicio. \\
GET  & \code{/api/mapas/estados}   & Los 32 estados con su clave INEGI. \\
GET  & \code{/api/mapas/talleres}  & Directorio DENUE filtrado por estado. Por defecto mecánica y llantas; el parámetro \code{categorias} agrega, por ejemplo, \code{carroceria} (hojalatería y pintura). \\
POST & \code{/api/mapas/solicitud} & Orquestador de agentes; responde 503 si no hay modelo configurado. \\
\bottomrule
\end{tabularx}
\caption{Rutas del servicio.}
\end{table}

%------------------------------------------------------------------------------
\section{Datos}

\subsection{Directorio de talleres (INEGI DENUE)}

Se usó la descarga masiva de DENUE para la CDMX, con corte 05/2026. Se conservan los
establecimientos de la rama 8111 del SCIAN M\'exico 2023 (reparación y mantenimiento de automóviles
y camiones). Las claves se verificaron contra el catálogo oficial y contra el propio archivo.
Se cargaron \textbf{15\,318} establecimientos, todos con coordenadas y 7\,768 con teléfono.

\begin{table}[h]
\centering
\begin{tabularx}{\textwidth}{@{}l Y r@{}}
\toprule
Categoría & Claves SCIAN & Talleres \\
\midrule
Mecánica   & 811111 a 811116 y 811119 (general, eléctrico, motor, transmisión, suspensión, alineación y balanceo, otras) & 9\,157 \\
Carrocería (opcional) & 811121 a 811123 y 811129 (hojalatería y pintura, tapicería, blindaje, cristales) & 2\,495 \\
Llantas    & 811191 (reparación menor de llantas) & 1\,786 \\
Servicios  & 811192 y 811199 (lavado y lubricado, otros) & 1\,880 \\
\bottomrule
\end{tabularx}
\caption{Establecimientos de la rama 8111 en la CDMX por categoría.}
\end{table}

\paragraph{Limitaciones.} DENUE no trae reseñas ni calificaciones. Tampoco distingue
``frenos'' ni ``reparación de motor'' como especialidades: los frenos caen dentro de la
mecánica general. Algunos registros tienen clasificación dudosa (por ejemplo, reparación de
autoestéreos dentro de mecánica) o no tienen nombre comercial. La base que viaja a producción
excluye teléfono, correo y sitio web.

\subsection{Red Nacional de Caminos (INEGI)}

La RNC 2025 se transformó en una base compacta nacional con \textbf{3\,343\,663} nodos y
\textbf{4\,349\,691} tramos (395 MB, fuera del repositorio). Cada tramo conserva longitud,
velocidad, sentido de circulación, estatus y peaje. Su licencia son los Términos de Libre Uso de INEGI.

\paragraph{Hallazgos que condicionan su uso.}
\begin{itemize}[nosep]
  \item El catálogo menciona la clase ``Taller Mecánico'', pero los datos no traen ningún registro de esa clase: la RNC no sirve como directorio de talleres.
  \item El campo de jurisdicción vale \code{N/A} en el 78\,\% de los tramos, por lo que no permite filtrar por estado.
  \item Las velocidades son fijas y no incluyen tráfico.
  \item \textbf{Supuesto sin verificar:} en los tramos de un solo sentido se circula de la unión inicial a la final.
\end{itemize}

%------------------------------------------------------------------------------
\section{Ranking de mecánicos}

\subsection{Cercanía}

La distancia en línea recta entre el usuario $(\varphi_1,\lambda_1)$ y un taller $(\varphi_2,\lambda_2)$ es
\begin{equation}
d = 2R\,\arcsin\!\sqrt{\sin^2\!\frac{\varphi_2-\varphi_1}{2}
      + \cos\varphi_1\cos\varphi_2\,\sin^2\!\frac{\lambda_2-\lambda_1}{2}},
\qquad R = 6371\ \text{km}.
\end{equation}
Cuando se implemente el tiempo de viaje con la RNC, sustituirá a $d$ en el desempate:
\begin{equation}
t = \sum_{e\,\in\,\text{ruta}} \frac{L_e}{v_e},
\end{equation}
con $L_e$ la longitud y $v_e$ la velocidad de cada tramo $e$.

\subsection{Experiencia con el problema}

Cada mecánico recibe un nivel discreto $E\in\{0,1,2,3\}$. La preferencia es
\begin{equation}
E=3 \;\succ\; E=2 \;\succ\; E=1 \;\succ\; E=0.
\end{equation}

\begin{center}
\begin{tabularx}{0.9\textwidth}{@{}c Y@{}}
\toprule
$E$ & Criterio \\
\midrule
0 & Sin evidencia. \\
1 & Declara la especialidad requerida. \\
2 & Al menos 1 trabajo verificado en esa especialidad con satisfacción $\ge 4$. \\
3 & Al menos 5 trabajos verificados así. \\
\bottomrule
\end{tabularx}
\end{center}

El nivel es discreto a propósito: con una variable continua nunca habría empates y los criterios siguientes no se usarían.

\subsection{Calificación ponderada por volumen}

Con $R$ la calificación del taller, $n$ su número de reseñas, $C$ el promedio global y $m$ el peso del prior,
\begin{equation}
\tilde{R} = \frac{n}{n+m}\,R + \frac{m}{n+m}\,C,
\qquad m = 50,\; C = 4{,}2 .
\end{equation}
Así un taller de 5 estrellas con una sola reseña no supera a uno de 4,5 con miles:

\begin{center}
\begin{tabular}{@{}l c c c@{}}
\toprule
Taller & $R$ & $n$ & $\tilde{R}$ \\
\midrule
A & 5,0 & 1      & 4,22 \\
B & 4,5 & 10\,000 & 4,50 \\
C & 4,8 & 20     & 4,37 \\
\bottomrule
\end{tabular}
\end{center}

\subsection{Clave lexicográfica y fases}

Los talleres que pasan las puertas de calidad (calificación mínima y número mínimo de reseñas) se
ordenan de forma ascendente por
\begin{equation}
\kappa(x) = \bigl(-E(x),\; -\tilde{R}(x),\; d(x)\bigr),
\end{equation}
comparando de izquierda a derecha: solo si coincide una componente se mira la siguiente. El flujo
``los 5 mejores y más cercanos'' usa la misma clave sin la componente $E$.

La búsqueda se hace en fases \emph{acumulativas} de \textbf{30, 60 y 100 talleres más cercanos}:
son cantidades de talleres, no distancias. Se pasa a la fase siguiente solo si en la actual nadie
pasa las puertas con $E\ge 1$. Si ninguna fase tiene un especialista, se informa y se ofrecen los
mejor calificados y cercanos. Con los datos reales de la CDMX, los 100 talleres más cercanos
caben en menos de 3,5 km en los puntos medidos, por lo que las fases importarán sobre todo en
estados poco densos; falta definir un tope de distancia.

\paragraph{Parámetros por revisar.} Los valores $m=50$ y las puertas (4,0 estrellas y 20 reseñas)
vienen de un contexto de reseñas de Google. Con reseñas propias, al arrancar casi ningún taller los
alcanzaría, así que son configurables y habrá que bajarlos.

%------------------------------------------------------------------------------
\section{Modelo de decisión lexicográfico}

Cada opción de diseño recibe un vector de puntajes $s=(s_1,\dots,s_6)\in\{1,2,3\}^6$, donde 3 es mejor,
uno por criterio. Dado un orden de prioridad $\pi$ de los criterios, la opción $a$ es preferida a $b$ si
existe una posición $k$ tal que
\begin{equation}
s_{\pi(j)}(a) = s_{\pi(j)}(b)\ \ \forall j<k
\quad\text{y}\quad
s_{\pi(k)}(a) > s_{\pi(k)}(b).
\end{equation}
El método no compensa: una ventaja en un criterio de mayor prioridad vale más que cualquier cantidad
de ventajas en los de menor prioridad. Por eso el orden es una decisión de valor del equipo.

\begin{table}[h]
\centering
\begin{tabularx}{\textwidth}{@{}l Y@{}}
\toprule
Criterio & Qué mide \\
\midrule
C1 & Cubre el flujo núcleo (problema, recomendación, cita). \\
C2 & Confianza de la recomendación. \\
C3 & Riesgo bajo: legal, privacidad, términos de uso y dinero. \\
C4 & Poco esfuerzo de implementación. \\
C5 & Bajo costo operativo. \\
C6 & Diferenciación y explicabilidad. \\
\bottomrule
\end{tabularx}
\caption{Criterios. El orden fijado por el equipo es C1 $>$ C3 $>$ C2 $>$ C4 $>$ C5 $>$ C6.}
\end{table}

\paragraph{Advertencia.} Los puntajes son juicios del equipo, no mediciones. Como contexto de
sensibilidad, para cada decisión se calculó en qué porcentaje de los 720 órdenes posibles de
criterios gana cada opción (columna \emph{robustez}).

\begin{table}[h]
\centering
\begin{tabularx}{\textwidth}{@{}l Y r@{}}
\toprule
Decisión & Opción que gana con el orden fijado & Robustez \\
\midrule
D1 Fuente de talleres & DENUE y mecánicos registrados & 75\,\% \\
D2 Reputación y experiencia & Reseñas propias verificadas (nivel $E$ y promedio bayesiano) & 75\,\% \\
D3 Prelista por cercanía & Los $k$ más cercanos: 30, 60 y 100 & 67\,\% \\
D4 Ranking final & Lexicográfico con puertas y promedio bayesiano & 50\,\% \\
D5 Agenda con el mecánico & Aprobación del mecánico, sin pago & 75\,\% \\
D10 Cercanía real & Grafo de la RNC en dos etapas & 47\,\% \\
D11 Modelo de lenguaje & Proveedor de pago (\textbf{pendiente} de confirmar con las claves disponibles) & 50\,\% \\
D12 Arquitectura & Todo en Vercel & 100\,\% \\
D13 Datos geográficos & SQLite embebida en la función (CDMX) & 67\,\% \\
\bottomrule
\end{tabularx}
\caption{Decisiones generadas por el modelo (ver \code{mapas/decisiones/DECISIONES.md}).}
\end{table}

%------------------------------------------------------------------------------
\section{Despliegue en Vercel}

Se usa Vercel Services (en beta según su documentación), que permite varios servicios en un mismo
proyecto. El archivo \code{vercel.json} de la raíz declara el servicio:

\begin{lstlisting}
{
  "services": {
    "mapas": { "root": "mapas/", "entrypoint": "main:app" }
  },
  "rewrites": [
    { "source": "/api/mapas/(.*)", "destination": { "service": "mapas" } }
  ]
}
\end{lstlisting}

\paragraph{Límites verificados en la documentación de Vercel.}
\begin{itemize}[nosep]
  \item Paquete de una función Python: 500 MB descomprimido (5 GB en beta).
  \item Duración máxima: 300 s en Hobby; 800 s en Pro. Esperar la respuesta de un modelo no cuenta como CPU activa.
  \item Memoria: 2 GB en Hobby; hasta 4 GB en Pro. Carga útil por petición o respuesta: 4,5 MB.
  \item El plan Hobby es solo para uso personal no comercial y no incluye funciones de equipo; Pro cuesta 20 USD por desarrollador al mes.
  \item Las funciones corren por defecto en una región de la costa este de EE.\,UU.; conviene acercarla a la base de datos.
\end{itemize}

\paragraph{Pasos.} (1) Importar el repositorio en Vercel. (2) Definir \code{LLM\_API\_KEY},
\code{LLM\_BASE\_URL} y \code{LLM\_MODEL} como variables de entorno del proyecto, nunca en el
repositorio. (3) Fijar la región de las funciones. (4) Desplegar una rama como \emph{preview} y
probar \code{/api/mapas/health}. El \code{vercel.json} sigue la documentación, pero
\textbf{no se ha probado en un despliegue real}.

%------------------------------------------------------------------------------
\section{Pruebas}

Hay 26 pruebas que no requieren red ni claves y que cubren el ordenamiento lexicográfico, los niveles
$E$ y el promedio bayesiano, las fases de búsqueda, el cargador y el filtro de estado del directorio,
y las rutas HTTP.

\begin{lstlisting}
cd mapas
pip install -r requirements-dev.txt
python -m pytest tests -q
\end{lstlisting}

No se han probado contra servicios reales: el modelo de lenguaje, la API de DENUE (se usó la descarga
masiva), ni el despliegue en Vercel.

%------------------------------------------------------------------------------
\section{Limitaciones y pendientes}

\begin{enumerate}[nosep]
  \item Proveedor de las claves del modelo y sus términos de uso en producción (D11).
  \item Cuenta propietaria y plan de Vercel (Hobby o Pro).
  \item Tiempo de viaje con la RNC y verificación del sentido de circulación de los tramos.
  \item Campo de estado en la tabla de mecánicos de la base de datos.
  \item Valores de $m$, $C$ y de las puertas de calidad para reseñas propias.
  \item Tope de distancia máxima para las fases de búsqueda en estados poco densos.
  \item Extender el directorio a otros estados (se carga un estado por archivo de DENUE).
  \item Reglas y tiempo disponible del reto.
\end{enumerate}

\section*{Fuentes}
\begin{itemize}[nosep]
  \item INEGI, Directorio Estadístico Nacional de Unidades Económicas (DENUE), descarga masiva, corte 05/2026.
  \item INEGI, SCIAN México 2023, catálogo de clases.
  \item INEGI, Red Nacional de Caminos 2025 (licencia: Términos de Libre Uso de la Información del INEGI).
  \item Vercel, documentación de Functions, runtime de Python, FastAPI y Services.
\end{itemize}

\end{document}
